\scanpage{89}

\margin{Transitivité $\swarrow$}

\deleted{[N.B. il suffit de le vérifier pour des ouverts \deleted{\unsure{affi}} affines quasi compacts \inserted{$U^\alpha$} recouvrant $X$, et pour tt $U^\alpha$ trouver un \uline{recouvrement} de $U^\alpha$ par des parties \deleted{constructibles}\inserted{loc. fermées} $U_i^\alpha$, tels que $f\mid |U_i^\alpha|$ soit \deleted{\unsure{aff}} morphique \ldots}

\margin{\deleted{Proposition}}

On appelle \uline{ensemble} \deleted{\unsure{algébrique}}\deleted{\unsure{arithmétique}} \uline{schématique} un \deleted{préschéma}\inserted{préschéma quasi-séparé quasi-compact}, considéré comme objet d'une catégorie dont les morphismes sont les applications constructibles.

\diagram{\tiny
\begin{tikzcd}[ampersand replacement=\&,column sep=6pt,row sep=1pt,
execute at end picture={\draw[thin] ([yshift=-8pt]\tikzcdmatrixname-1-1.south east) -- ++(6pt,5pt) -- ++(-4pt,-9pt) -- ([yshift=-6pt]\tikzcdmatrixname-1-2.south west) -- ++(-4pt,-9pt) -- ([yshift=-6pt]\tikzcdmatrixname-1-2.south) -- ++(-3pt,-10pt) -- ([yshift=-6pt]\tikzcdmatrixname-1-3.south west) -- ++(-3pt,-10pt) -- ([yshift=-9pt]\tikzcdmatrixname-1-4.south west);}]
{(\text{Schs }\deleted{\unsure{\text{réd}}})}\arrow[r]
\arrow[rr,bend right=23,"{\deleted{\text{fidèle}}}"']
\arrow[rrr,bend right=30,"{\deleted{\text{fidèle}}}"']
\& {(\text{Schs})}\arrow[r,"{\unsure{\text{n'd}|\text{réd}|\text{red}}\ \text{fidèle}}"']
\arrow[rr,bend right=21,"{\deleted{\text{fidèle}}}"']
\& {(\text{ens Schs})}\arrow[r,"{\text{fidèle}}"']
\& {(\text{Sch. \unsure{atomes|atoms|atom.}})}\arrow[r,"{\text{pl. fid.}}"']
\& {\text{\uline{Hom}}\,((\text{Corps}),\text{Ens})}
\end{tikzcd}}

\margin{\fbox{\parbox{1.8cm}{N.B. Tous les foncteurs écrits \unsure{ici|ci-dessus} commutent aux $\underset{\leftarrow}{\lim}$ finies}}}

On s'intéresse \unsure{donc|données} aux notions sur les ens. schématiques, qui sont invariantes par \unsure{isom|isom.} d'ens schématiques (i.e. isom. constructibles).

Notons \deleted{\unsure{y|$X$}} comme \unsure{Rem. que|lemme que}

\diagram{\begin{tikzcd}[ampersand replacement=\&,row sep=2pt]
X\arrow[r]\&Y\\
X_i\&Y_i\\
X'_j\&\\
X_i\cap X'_j\arrow[r]\&Y_{ij}
\end{tikzcd}}

\uline{Lemme} Soit $X$ \deleted{\unsure{pré}}préschéma quasi-compact, quasi-séparé. \deleted{Alors il existe une} $(Z_i)$ une famille \deleted{\unsure{de|finie|disjointe}} de parties constructibles recouvrant $X$. Alors il existe une partition de $X$ en des parties loc. fermées constructibles affines, plus fine que \unsure{un|une} \unsure{rec. fini|recouvrement} \uline{\unsure{donné}} de $X$ par des ens. constructibles

\uline{Proposition} Supposons $X$ quasi-séparé et quasi-compact, $f:|X|\to|Y|$ donné. Conditions équivalentes

a) $f$ est constructible

\unsure{aiii)|a ii)|(ii)} $\exists$ une partition \inserted{finie} de $X$ en des parties loc. \uline{fermées constructibles} \inserted{affines} $Z_i$, telles que $f\mid |Z_i|$ soit morphique

\deleted{\unsure{(iii)|bii)} Il existe une suite croissante d'ouverts quasi-compacts $U_i$ couvrant $X$, telle que $f\mid |U_i-U_{i-1}|$ soit \unsure{morphique}.}
% archive sheet number: 88
